% !TEX root = ../access.tex
%% ============================================================
%% SECTION 4 — RESULTS AND DISCUSSION
%% Showcases a full-width (two-column-spanning) figure and table.
%% ============================================================
\section{Results and Discussion}
\label{sec:results}

Present and discuss your results. Reference the wide table as
Table~\ref{tab:wide} and the wide figure as Fig.~\ref{fig:wide}.

%% ---- Full-width table (spans both columns) ---------------
% Use table* + tabular*{\textwidth} + \extracolsep{\fill} so the
% table fills the full text width across both columns.
\begin{table*}[!tb]
\caption{Full-Width Comparison Table Example}
\label{tab:wide}
\setlength{\tabcolsep}{4pt}
\begin{tabular*}{\textwidth}{@{\extracolsep{\fill}}lccccc@{}}
\hline
\textbf{Method} & \textbf{Metric A} & \textbf{Metric B} & \textbf{Metric C}
& \textbf{Metric D} & \textbf{Metric E} \\
\hline
Baseline         & 00.0 & 00.0 & 00.0 & 00.0 & 00.0 \\
Competitor       & 00.0 & 00.0 & 00.0 & 00.0 & 00.0 \\
\textbf{Ours}    & \textbf{00.0} & \textbf{00.0} & \textbf{00.0}
                 & \textbf{00.0} & \textbf{00.0} \\
\hline
\end{tabular*}
\end{table*}

%% ---- Full-width figure (spans both columns) --------------
\begin{figure*}[!t]
\centering
% \includegraphics[width=\textwidth]{wide_figure.png}
\caption{Full-width figure caption. Use \texttt{figure*} +
\texttt{width=\textbackslash textwidth} for a figure that spans both columns.}
\label{fig:wide}
\end{figure*}

Discuss what the numbers mean, compare against baselines, and note
limitations.
